\section{Adversário Adaptativo (conhece a defesa)}\label{sec:adaptive}
Avaliamos um atacante onisciente que conhece a regra de corroboração, o limiar $T$ e as
atualizações honestas da rodada. Ele (A) concentra o envenenamento em $m$ tensores --- revertendo
os demais tensores sensíveis para sua própria atualização benigna, para que não sinalizem --- e (B)
dessincroniza os clientes maliciosos com ruído independente, de modo que o cosseno mútuo caia; ele
varre $(m,\nu)$ para encontrar o backdoor mais forte que mantém a contagem de corroboração abaixo de
$T$. A hipótese sob teste (H4) é que há uma tensão fundamental: reduzir $m$ para escapar da contagem
também reduz a fração de parâmetros do backdoor, enfraquecendo o ataque.
\textbf{Resultado.} O atacante adaptativo vence a corroboração (Tabela~\ref{tab:adapt}, sementes 42--47):
a ASR final é $62{,}5\%$ contra a \emph{corr} e $53{,}5\%$ contra a \emph{corr-soft}, ante $83{,}4\%$ sem
defesa com o mesmo ataque concentrado. A \emph{client} resiste mais ($25{,}0\%$ em média), mas é instável
(de $6{,}0$ a $58{,}5\%$ entre sementes) e perde $8{,}4$\,pp de acurácia. A hipótese H4 se confirma só em
parte: o atacante raramente consegue deixar \emph{todos} os maliciosos abaixo de $T$ na mesma rodada
($1$--$3\%$ das rodadas), mas não precisa disso. Em $37\%$ dos pares (atacante, rodada) contra a
\emph{corr-soft}, um malicioso fica abaixo de $T$ e entra na agregação sem ser confirmado; o veneno passa
aos poucos, por dois ou três atacantes a cada rodada, e o backdoor se acumula.

\textbf{O custo imposto ao atacante.} A defesa não impede o ataque, mas o encarece. Contra a
\emph{corr-soft}, o backdoor leva o dobro do tempo para pegar: chega a $50\%$ na rodada $165$, ou seja,
$85$ rodadas depois do início do ataque, contra $43$ rodadas sem defesa; a ASR média durante o ataque
cai de $62{,}1$ para $37{,}7\%$; a final, de $83{,}4$ para $53{,}5\%$ ($-30$\,pp). Contra a \emph{client},
o ataque só passa de $20\%$ por volta da rodada $167$ e, em 5 das 6 sementes, nunca chega a $50\%$.

\begin{table}[t]\centering\small
\caption{Atacante adaptativo ciente da corroboração, cenário fiel do LaySelFL, sementes 42--47 (médias).
O ataque começa na rodada 80. Valores de ASR e acurácia em \%.}
\label{tab:adapt}
\begin{tabular}{lccccc}
\toprule
Agregação & ASR final & ASR média no ataque & ASR $\geq 50\%$ na rodada & Acurácia \\
\midrule
Sem defesa (ataque concentrado) & $83{,}4$ & $62{,}1$ & $123$ & $76{,}8$ \\
corr & $62{,}5$ & $43{,}3$ & $155$ & $74{,}8$ \\
\textbf{corr-soft} & $53{,}5$ & $37{,}7$ & $165$ & $74{,}5$ \\
client & $25{,}0$ & $15{,}5$ & $180$ (5 de 6 nunca) & $68{,}4$ \\
\bottomrule
\end{tabular}
\end{table}

\textbf{Respostas testadas.} Aplicar uma agregação robusta depois da corroboração não resolveu, contra um
atacante que também conhece essa agregação (sementes 42--45): com corte de norma por tensor a ASR final
fica em $60{,}2\%$ e com média aparada em $72{,}5\%$. Um atacante intermediário, que conhece a regra da
contagem mas não simula a agregação, chega a $62{,}8\%$ contra o corte de norma. Duas variantes que atacam
diretamente o vazamento abaixo de $T$ --- peso gradual em função da contagem e $T$ sorteado em segredo a
cada rodada --- estão em avaliação. \todo{inserir resultado da Fase 20.}
